\documentclass[10pt,a4paper]{amsart}
\usepackage{amsmath,amssymb,amsthm,graphicx}
\usepackage{verbatim}
\newenvironment{reference}{\comment}{\endcomment}
%\newenvironment{reference}{}{}
\newenvironment{slogan}{\comment}{\endcomment}
\newenvironment{history}{\comment}{\endcomment}

% For commutative diagrams we use Xy-pic
\usepackage[all]{xy}

% We use 2cell for 2-commutative diagrams.
\xyoption{2cell}
% Theorem environments.
%
\theoremstyle{plain}
\newtheorem{theorem}[subsection]{Theorem}
\newtheorem{proposition}[subsection]{Proposition}
\newtheorem{lemma}[subsection]{Lemma}

\theoremstyle{definition}
\newtheorem{definition}[subsection]{Definition}
\newtheorem{example}[subsection]{Example}
\newtheorem{exercise}[subsection]{Exercise}
\newtheorem{situation}[subsection]{Situation}

\theoremstyle{remark}
\newtheorem{remark}[subsection]{Remark}
\newtheorem{remarks}[subsection]{Remarks}

\numberwithin{equation}{subsection}

% Macros
%
\def\lim{\mathop{\mathrm{lim}}\nolimits}
\def\colim{\mathop{\mathrm{colim}}\nolimits}
\def\Spec{\mathop{\mathrm{Spec}}}
\def\Hom{\mathop{\mathrm{Hom}}\nolimits}
\def\Ext{\mathop{\mathrm{Ext}}\nolimits}
\def\SheafHom{\mathop{\mathcal{H}\!\mathit{om}}\nolimits}
\def\SheafExt{\mathop{\mathcal{E}\!\mathit{xt}}\nolimits}
\def\Sch{\mathit{Sch}}
\def\Mor{\mathop{\mathrm{Mor}}\nolimits}
\def\Ob{\mathop{\mathrm{Ob}}\nolimits}
\def\Sh{\mathop{\mathit{Sh}}\nolimits}
\def\NL{\mathop{N\!L}\nolimits}
\def\CH{\mathop{\mathrm{CH}}\nolimits}
\def\proetale{{pro\text{-}\acute{e}tale}}
\def\etale{{\acute{e}tale}}
\def\QCoh{\mathit{QCoh}}
\def\Ker{\mathop{\mathrm{Ker}}}
\def\Im{\mathop{\mathrm{Im}}}
\def\Coker{\mathop{\mathrm{Coker}}}
\def\Coim{\mathop{\mathrm{Coim}}}

% Boxtimes
%
\DeclareMathSymbol{\boxtimes}{\mathbin}{AMSa}{"02}

%
% Macros for moduli stacks/spaces
%
\def\QCohstack{\mathcal{QC}\!\mathit{oh}}
\def\Cohstack{\mathcal{C}\!\mathit{oh}}
\def\Spacesstack{\mathcal{S}\!\mathit{paces}}
\def\Quotfunctor{\mathrm{Quot}}
\def\Hilbfunctor{\mathrm{Hilb}}
\def\Curvesstack{\mathcal{C}\!\mathit{urves}}
\def\Polarizedstack{\mathcal{P}\!\mathit{olarized}}
\def\Complexesstack{\mathcal{C}\!\mathit{omplexes}}
% \Pic is the operator that assigns to X its picard group, usage \Pic(X)
% \Picardstack_{X/B} denotes the Picard stack of X over B
% \Picardfunctor_{X/B} denotes the Picard functor of X over B
\def\Pic{\mathop{\mathrm{Pic}}\nolimits}
\def\Picardstack{\mathcal{P}\!\mathit{ic}}
\def\Picardfunctor{\mathrm{Pic}}
\def\Deformationcategory{\mathcal{D}\!\mathit{ef}}
\usepackage{cleveref}
\crefname{originallemma}{original lemma}{original lemmas}
\Crefname{originallemma}{Original lemma}{Original lemmas}
\makeatletter
\newlabel{lemma-cartesian-cocartesian}{{08N4}{source}{}{sourcepointer-1}{}}
\newlabel{lemma-cartesian-cocartesian@cref}{{[originallemma][0][]08N4}{[1][0][]source}{}{}{}}
\makeatother
\title{A criterion for exactness through a fibre product}
\author{The Stacks Project Authors}
\date{stacks-project:homology.tex:local-snapshot-20261012; GNU FDL 1.2 or later}
\begin{document}
\maketitle
\noindent\textit{Editorial note.} All delivered mathematical text is the author's own, copied from the pinned Stacks Project chapter file homology.tex; no statement or proof step was rewritten, reordered or removed.\par\medskip
\noindent\textit{Editorial note.} The proof is the author's complete argument in both directions; every step is inside the counted span, the cancellations are justified by naming the monomorphism or epimorphism being cancelled, and no sentence of the proof is a placeholder.\par\medskip
\noindent\textit{Editorial note.} Every labelled reference occurring inside the counted spans is carried in assembly.external\_references as a non-load-bearing navigation pointer (same-chapter ones at their definition site, cross-chapter ones at the invocation inside the counted span), so no load-bearing mathematics is replaced by a stub. No bibliography citation occurs inside the counted spans.\par\medskip
\noindent\textit{Editorial note.} The same-name check for this card scanned every staged Stacks chapter of this repository -- all sources/stacks-*/pinned-original-source.tex files, the card's own chapter excluded -- and the labels found are recorded in dedup.same\_name\_cross\_chapter on the card.\par\medskip
\noindent\textit{Editorial note.} The preamble is the chapter's own native preamble excerpt already used for this source: the comment-type reference, slogan and history environments (native preamble lines 8-13), the Xy-pic setup (native preamble lines 14-18) and the native theorem declarations with the mathematical macros (native preamble lines 154-217, the actual last line of the file).\par\medskip
\noindent\textit{Allowed standard prerequisites (editorial).}\begin{itemize}
\item Cartesian and cocartesian squares. Same chapter, Lemma lemma-cartesian-cocartesian. Supplies that the projection from a fibre product is an epimorphism when the other structure map is one, which is how the epimorphism demanded by condition (2) is produced.
\item Coimage, image and the canonical factorisation. Same chapter, the definitions of the coimage and the image of a morphism in an abelian category. Supplies the canonical projection onto the coimage, the canonical injection of the image into the kernel of the next map, and their compatibility.
\end{itemize}\medskip
\ifdefined\hypertarget\hypertarget{sourcepointer-1}{}\fi
\noindent\textit{Source reference.} Original lemma 08N4: sources/stacks-homology/pinned-original-source.tex, source lines 1104--1104. Navigation pointer invoked at proof line 1148; the counted span contains the pointer only and the cited result is not reproduced inside it, so the pointer is registered as a non-load-bearing navigation stub in the established form for Stacks material. No bibliography citation is involved.\par\medskip
\begin{lemma}
\label{lemma-check-exactness-fibre-product}
Let $\mathcal{A}$ be an abelian category. Let $f:x\to y$ and $g:y\to z$ 
be morphisms with $g\circ f=0$. Then, the following statements are equivalent:
\begin{enumerate}
\item The sequence $x\overset{f}\to y\overset{g}\to z$ is exact.
\item For every $h:w\to y$ with $g\circ h=0$ there exist an object $v$, 
an epimorphism $k:v\to w$ and a morphism $l:v\to x$ with $h\circ k=f\circ l$.
\end{enumerate}
\end{lemma}

\begin{proof}
Let $i:\Ker(g)\to y$ be the canonical injection. Let 
$p:x\to\Coim(f)$ be the canonical projection. Let 
$j:\Im(f)\to\Ker(g)$ be the canonical injection.

\medskip\noindent
Suppose (1) holds. Let $h:w\to y$ with $g\circ h=0$. There exists 
$c:w\to\Ker(g)$ with $i\circ c=h$. 
Let $v=x\times_{\Ker(g)}w$ with canonical projections 
$k:v\to w$ and $l:v\to x$, so that $c\circ k=j\circ p\circ l$. 
Then, $h\circ k=i\circ c\circ k=i\circ j\circ p\circ l=f\circ l$. 
As $j\circ p$ is an epimorphism by hypothesis, $k$ is an 
epimorphism by Lemma \ref{lemma-cartesian-cocartesian}. This implies (2).

\medskip\noindent
Suppose (2) holds. Then, $g\circ i=0$. So, there are an object 
$w$, an epimorphism $k:w\to\Ker(g)$ and a morphism 
$l:w\to x$ with $f\circ l=i\circ k$. It follows 
$i\circ j\circ p\circ l=f\circ l=i\circ k$. Since $i$ is a 
monomorphism we see that $j\circ p\circ l=k$ is an epimorphism. 
So, $j$ is an epimorphisms and thus an isomorphism. This implies (1).
\end{proof}

\end{document}
